% Reviewed translation: PDF pages 245--250 / printed pages 231--236.
% The original scan is authoritative; mathematical tokens were checked against it.
\MathBookSourcePage{245}
\subsection{$l$ 次四边形有限元的超收敛}
\label{subsec:incompressible-quadrilateral-superconvergence}

本小节再次研究表达式
\[
\inf_{\phi_h\in\Phi_h}
\sup_{\mu_h\in\Theta_h}
\frac{|(\operatorname{curl}(\phi-\phi_h),
\operatorname{curl}\mu_h)|}{\|\mu_h\|_{0,\Omega}},
\]
并证明, 若剖分 $\mathcal T_h$ 具有有利的构形,
则对 $\phi\in H^{l+2}(\Omega)$, 它是 $h^l$ 阶的.
这主要通过把 $\phi_h$ 选成 $\phi$ 的适当插值来实现.

首先假设 $\mathcal T_h$ 完全由凸四边形组成,
并采用 Section~\ref{sec:appendix-quadrilateral-finite-elements}
中关于四边形的记号和逼近结果.
固定整数 $l\geq1$, 令
\[
0\leq\alpha_0<\alpha_1<\cdots<\alpha_l\leq1
\]
为区间 $[0,1]$ 的 $l+1$ 个不同点.
记 $I\in\mathcal L(C^0(\widehat\kappa);Q_l)$
为参考单位正方形 $\widehat\kappa$ 上节点
\[
\{(\alpha_i,\alpha_j);\ 0\leq i,j\leq l\}
\]
处的标准插值算子.
下列技术结果给出差 $\phi-I\phi$ 的第一个逼近.

\begin{Lemma}\label{lem:incompressible-superconvergence-interpolation-remainder}
令 $R\in\mathcal L(H^{l+2}(\widehat\kappa);P_{l+1,l+1})$ 定义为
\[
R(\phi)=\frac1{(l+1)!}
\left\{
\frac{\partial^{l+1}\phi}{\partial x_1^{l+1}}(0,x_2)
\prod_{i=0}^l(x_1-\alpha_i)
+\frac{\partial^{l+1}\phi}{\partial x_2^{l+1}}(x_1,0)
\prod_{j=0}^l(x_2-\alpha_j)
\right\}.
\]
则对每个整数 $m$, $0\leq m\leq l+2$,
存在只依赖于 $\widehat\kappa$, $I$, $R$, $m$ 和 $l$ 的正常数 $C$, 使
\begin{equation}\label{eq:incompressible-superconvergence-interpolation-remainder}
\|\phi-I\phi-R(\phi)\|_{m,\widehat\kappa}
\leq C|\phi|_{l+2,\widehat\kappa}
\qquad\forall\phi\in H^{l+2}(\widehat\kappa).
\end{equation}
\end{Lemma}

\begin{Proof}
先设 $\phi\in C^{l+1}(\widehat\kappa)$.
回顾插值算子可写成
\[
I\phi(x_1,x_2)
=\sum_{i=0}^l\sum_{j=0}^l
L_i(x_1)L_j(x_2)\phi(\alpha_i,\alpha_j),
\]
其中 $L_i\in P_l$ 且 $L_i(\alpha_j)=\delta_{ij}$.
熟知的插值误差余项公式分别沿两个坐标方向应用,
可把 $\phi-I\phi$ 写成由
$\prod_{i=0}^l(x_1-\alpha_i)$ 和
$\prod_{j=0}^l(x_2-\alpha_j)$ 相乘的两个
$(l+1)$ 阶导数项之和.
将它与 $R(\phi)$ 比较, 并使用
$\sum_{i=0}^lL_i(x)=1$, 容易看出映射
$\phi\mapsto\phi-I\phi-R(\phi)$ 在 $P_{l+1,l+1}$ 上为零.
该映射属于
$\mathcal L(H^{l+2}(\widehat\kappa);H^m(\widehat\kappa))$,
因此 \eqref{eq:incompressible-superconvergence-interpolation-remainder}
由 Corollary~\ref{cor:appendix-abstract-interpolation-reference} 得到.
\end{Proof}

下一个结果通过简单重排各项得到.

\MathBookSourcePage{246}
\begin{Lemma}\label{lem:incompressible-superconvergence-mapped-gradient-identity}
对所有 $\phi,\mu\in H^1(\kappa)$, 有
\begin{equation}\label{eq:incompressible-superconvergence-mapped-gradient-identity}
\begin{split}
\int_\kappa\operatorname{grad}\phi\cdot
\operatorname{grad}\mu\,dx
={}&\int_{\widehat\kappa}\frac1{J_\kappa}
\biggl[
\left\|\frac{\partial F_\kappa}{\partial\widehat x_2}\right\|^2
\frac{\partial\widehat\phi}{\partial\widehat x_1}
\frac{\partial\widehat\mu}{\partial\widehat x_1}
+\left\|\frac{\partial F_\kappa}{\partial\widehat x_1}\right\|^2
\frac{\partial\widehat\phi}{\partial\widehat x_2}
\frac{\partial\widehat\mu}{\partial\widehat x_2}
\biggr]d\widehat x\\
&-\int_{\widehat\kappa}\frac1{J_\kappa}
\biggl[
\frac{\partial F_\kappa}{\partial\widehat x_1}\cdot
\frac{\partial F_\kappa}{\partial\widehat x_2}
\left(
\frac{\partial\widehat\phi}{\partial\widehat x_1}
\frac{\partial\widehat\mu}{\partial\widehat x_2}
+\frac{\partial\widehat\phi}{\partial\widehat x_2}
\frac{\partial\widehat\mu}{\partial\widehat x_1}
\right)\biggr]d\widehat x,
\end{split}
\end{equation}
其中 $\|\cdot\|$ 表示 $\mathbb R^2$ 中的 Euclidean 范数,
$J_\kappa$ 是双线性映射 $F_\kappa$ 的 Jacobian.
\end{Lemma}

为恰当地估计
\eqref{eq:incompressible-superconvergence-mapped-gradient-identity}
右端, 需要把含 $J_\kappa$ 和
$\partial F_\kappa/\partial\widehat x_i$ 的因子移出积分.
因为它们不是常数, 必须假设其导数 "很小".
具体作如下假设: 存在与 $h$ 和 $\kappa$ 无关的常数 $C>0$, 使
\begin{equation}\label{eq:incompressible-superconvergence-almost-parallelogram}
\left|\frac{\partial^2F_\kappa}
{\partial\widehat x_1\partial\widehat x_2}\right|
\leq Ch_\kappa^2
\qquad\forall\kappa\in\mathcal T_h.
\end{equation}
当 $\kappa$ 是平行四边形时, 上式左端为零;
事实上, 它只在 $\kappa$ 几乎为平行四边形时成立.
容易证明, 若
\eqref{eq:incompressible-superconvergence-almost-parallelogram}
成立, 则全半范数 $|\cdot|_{k,\kappa}$
与方向半范数 $[\cdot]_{k,\kappa}$
(参见 \eqref{eq:appendix-directional-seminorm})
具有同阶上界.

\begin{Lemma}\label{lem:incompressible-superconvergence-map-bounds}
令凸四边形 $\kappa$ 满足
\eqref{eq:incompressible-superconvergence-almost-parallelogram}.
则对每个整数 $k\geq1$, 存在与 $h$ 和 $\kappa$ 无关的常数 $C>0$, 使
\begin{equation}\label{eq:incompressible-superconvergence-directional-bound}
|\phi|_{k,\kappa}
\leq C\left(\sum_{i=1}^k[\phi]_{i,\kappa}^2\right)^{1/2}
\qquad\forall\phi\in H^k(\kappa).
\end{equation}
此外,
\begin{equation}\label{eq:incompressible-superconvergence-coefficient-derivative}
\left\|
\frac{\partial}{\partial\widehat x_k}
\left[
\frac1{J_\kappa}
\frac{\partial F_\kappa}{\partial\widehat x_i}
\cdot\frac{\partial F_\kappa}{\partial\widehat x_j}
\right]\right\|_{0,\infty,\widehat\kappa}
\leq C\sigma h_\kappa(1+\sigma),
\end{equation}
其中 $1\leq i,j,k\leq2$.
\end{Lemma}

\MathBookSourcePage{247}
若剖分 $\mathcal T_h$ 正则,
\eqref{eq:incompressible-superconvergence-coefficient-derivative}
蕴含
\[
X_\kappa(\widehat x_1,\widehat x_2)
=X_\kappa(1/2,1/2)+R_\kappa,
\quad
X_\kappa=\frac1{J_\kappa}
\frac{\partial F_\kappa}{\partial\widehat x_i}
\cdot\frac{\partial F_\kappa}{\partial\widehat x_j},
\]
且余项满足 $|R_\kappa|\leq C\sigma^2h_\kappa$.
由于我们特别关注恢复一个 $h$ 次幂, 可把这个余项计入高阶项.
因此, 由
\eqref{eq:incompressible-superconvergence-mapped-gradient-identity}
和 \eqref{eq:incompressible-superconvergence-coefficient-derivative},
只需研究四类项
\begin{equation}\label{eq:incompressible-superconvergence-four-terms}
\int_{\widehat\kappa}
\frac{\partial(\widehat\phi-\widehat\phi_h)}
{\partial\widehat x_i}
\frac{\partial\widehat\mu_h}{\partial\widehat x_j}
\,d\widehat x,
\qquad1\leq i,j\leq2.
\end{equation}
若取 $\phi_h=I_\kappa\phi$, 其中 $I_\kappa$
是在 $\kappa$ 上与 $I$ 对应的插值算子,
则 Lemma~\ref{lem:incompressible-superconvergence-interpolation-remainder}
表明, 当 $i=j$ 时,
\eqref{eq:incompressible-superconvergence-four-terms}
特别含有形如
\[
\int_0^1w_i(x_i)
\frac{\partial\widehat\mu_h}{\partial\widehat x_i}
\,d\widehat x_i,
\qquad
w_i(x_i)=\prod_{p=0}^l(x_i-\alpha_p)
\]
的因子. 若 $\alpha_p$ 是高精度求积公式的节点,
该积分可能为零. 这促使我们选择 $[0,1]$
上 Gauss--Lobatto 求积公式的 $l+1$ 个节点;
回顾该公式对 $2l-1$ 次多项式精确.

\begin{Lemma}\label{lem:incompressible-superconvergence-diagonal-terms}
令 $\kappa$ 如
Lemma~\ref{lem:incompressible-superconvergence-map-bounds} 中所述,
$\{\alpha_i\}_{i=0}^l$ 为 Gauss--Lobatto 求积公式的 $l+1$ 个节点.
则存在与 $h$ 和 $\kappa$ 无关的常数 $C>0$, 使
\begin{equation}\label{eq:incompressible-superconvergence-diagonal-terms}
\left|
\int_{\widehat\kappa}
\frac{\partial(\widehat\phi-I\widehat\phi)}
{\partial\widehat x_i}
\frac{\partial\widehat\mu_h}{\partial\widehat x_i}
\,d\widehat x
\right|
\leq C\sigma^2h_\kappa^l
\|\phi\|_{l+2,\kappa}\|\mu_h\|_{0,\kappa},
\quad i=1,2.
\end{equation}
\end{Lemma}

\MathBookSourcePage{248}
\begin{Proof}
如上所述,
Lemma~\ref{lem:incompressible-superconvergence-interpolation-remainder}
把左端分成 $R(\widehat\phi)$ 所给主项与余项 $E_i$,
其中由
\eqref{eq:incompressible-superconvergence-directional-bound}
和附录尺度估计,
\[
|E_i|
\leq C_1|\widehat\phi|_{l+2,\widehat\kappa}
|\widehat\mu_h|_{1,\widehat\kappa}
\leq C_3\sigma^2h_\kappa^l
\|\phi\|_{l+2,\kappa}\|\mu_h\|_{0,\kappa}.
\]
对主项作一次分部积分, 得
\[
\int_0^1w_i(x_i)
\frac{\partial f_i}{\partial x_i}\,dx_i=0.
\]
这里 $w_i/f_i$ 在 $\widehat\kappa$ 的边界上为零,
而被积函数 $w_i\partial^2f_i/\partial x_i^2$
属于 $Q_{2l-1}$, 并在节点 $\{\alpha_i\}$ 上为零.
这证明 \eqref{eq:incompressible-superconvergence-diagonal-terms}.
\end{Proof}

该 Lemma 处理了
\eqref{eq:incompressible-superconvergence-four-terms}
中前两项. 但当 $i\ne j$ 时的第三、第四项不易处理.
从 $R(\phi)$ 可见, 上述论证会产生 $Q_{2l}$ 中多项式的积分,
而 Gauss--Lobatto 公式只对 $2l-1$ 次多项式精确,
所以节点的特殊选择在这里无效.

仔细考察交叉项, 为消去 $\mu_h$ 的导数,
对 $\widehat x_j$ 分部积分.
再次应用 Lemma~\ref{lem:incompressible-superconvergence-interpolation-remainder}
并利用 $i\ne j$ 时
$\partial^2R(\phi)/(\partial\widehat x_i\partial\widehat x_j)=0$, 得
\[
\int_{\widehat\kappa}
\frac{\partial(\widehat\phi-I\widehat\phi)}
{\partial\widehat x_i}
\frac{\partial\widehat\mu_h}{\partial\widehat x_j}
\,d\widehat x
=\int_{\partial\widehat\kappa}
\frac{\partial(\widehat\phi-I\widehat\phi)}
{\partial\widehat x_i}\widehat\mu_h n_j\,ds+E_i,
\]
其中
\begin{equation}\label{eq:incompressible-superconvergence-cross-remainder}
|E_i|\leq C_1\sigma^2h_\kappa^l
\|\phi\|_{l+2,\kappa}\|\mu_h\|_{0,\kappa}.
\end{equation}

\MathBookSourcePage{249}
但对 $i\ne j$,
$[\partial(\phi-I_\kappa\phi)/\partial x_i]n_j$
在单元边界间连续而符号相反, 并在 $\Gamma$ 上为零.
单独估计一个单元只能得到
\begin{equation}\label{eq:incompressible-superconvergence-crude-cross-bound}
\left|
\int_{\partial\kappa}
\frac{\partial(\phi-I_\kappa\phi)}{\partial x_i}
\mu_h n_j\,ds\right|
\leq C\|\phi\|_{l+1,\kappa}\|\mu_h\|_{0,\kappa}
\leq C\sigma h_\kappa^{l-1}
\|\phi\|_{l+1,\kappa}\|\mu_h\|_{0,\kappa}.
\end{equation}
因此为改进这个估计, 必须先对 $\mathcal T_h$ 的所有单元求和,
再估计它们在给定内部边界上的贡献.
只有当相邻两个单元产生的因子
\[
X_\kappa
=\frac1{J_\kappa}
\frac{\partial F_\kappa}{\partial\widehat x_1}
\cdot\frac{\partial F_\kappa}{\partial\widehat x_2}
\]
之差 "很小" 时, 才能得到更好估计.

令 $\kappa_1,\kappa_2$ 为两个相邻单元.
由 \eqref{eq:incompressible-superconvergence-coefficient-derivative},
只需估计
$X_{\kappa_1}(a_1,b_1)-X_{\kappa_2}(a_2,b_2)$,
其中 $(a_i,b_i)$ 是 $\kappa_i$ 上任意一点.
最简单的选择是在公共边 $\kappa'$ 上取一点 $(a,b)$,
并置 $(a_1,b_1)=(a_2,b_2)=(a,b)$.
令 $i$ 是使 $\partial F_{\kappa}/\partial\widehat x_i$
跨过公共边连续的指标. 直接计算得
\[
|J_{\kappa_1}J_{\kappa_2}
(X_{\kappa_1}-X_{\kappa_2})(a,b)|
=\left\|\frac{\partial F}{\partial\widehat x_i}(a,b)\right\|^2
\left|
\frac{\partial F_{\kappa_1}}{\partial\widehat x_j}
\times
\frac{\partial F_{\kappa_2}}{\partial\widehat x_j}(a,b)
\right|,
\quad j\ne i.
\]
取 $(a,b)$ 为 $\kappa'$ 的一个端点,
令 $L_1,L_2$ 分别表示在该端点与 $\kappa'$ 相交的
$\kappa_1,\kappa_2$ 的另一条边. 则
\begin{equation}\label{eq:incompressible-superconvergence-adjacent-geometry}
|J_{\kappa_1}J_{\kappa_2}
(X_{\kappa_1}-X_{\kappa_2})(a,b)|
=\operatorname{meas}^2(\kappa')
\operatorname{meas}(L_1)\operatorname{meas}(L_2)
|\mathbf t_1\cdot\mathbf n_2|,
\end{equation}
其中 $\mathbf t_1$ 是 $L_1$ 的单位方向向量,
$\mathbf n_2$ 是 $L_2$ 的单位法向量.

\MathBookSourcePage{250}
因此, 若 $|\mathbf t_1\cdot\mathbf n_2|$ 也是 $O(h)$,
则 $X_{\kappa_1}-X_{\kappa_2}$ 为 $O(h)$;
换言之, $L_1,L_2$ 几乎平行.
采用如下假设: 存在与 $h$ 无关的常数 $C>0$, 使
\begin{equation}\label{eq:incompressible-superconvergence-nearly-parallel}
|\mathbf t_1\cdot\mathbf n_2|\leq Ch
\end{equation}
对 $\mathcal T_h$ 中所有这样的相邻线段对 $L_1,L_2$ 成立.
该假设意味着, 同时满足
\eqref{eq:incompressible-superconvergence-almost-parallelogram}
和 \eqref{eq:incompressible-superconvergence-nearly-parallel}
的网格可由两束平行直线作小扰动得到.
这当然是很严格的网格条件, 只有很少的区域 $\Omega$
容易接受这种剖分. 尽管如此, 它仍给出本小节开头宣布的结果.

\begin{Lemma}\label{lem:incompressible-superconvergence-cross-terms}
令 $\Omega$ 是有界多边形,
$\mathcal T_h$ 是由满足
\eqref{eq:incompressible-superconvergence-almost-parallelogram}
和 \eqref{eq:incompressible-superconvergence-nearly-parallel}
的凸四边形构成的正则剖分.
则存在与 $h$ 无关的常数 $C>0$, 使当 $i\ne j$ 时
\begin{equation}\label{eq:incompressible-superconvergence-cross-terms}
\sum_{\kappa\in\mathcal T_h}
\left|
\int_{\widehat\kappa}
\frac{\partial(\widehat\phi-I\widehat\phi)}
{\partial\widehat x_i}
\frac{\partial\widehat\mu_h}{\partial\widehat x_j}
\,d\widehat x
\right|
\leq Ch^l\|\phi\|_{l+2,\Omega}\|\mu_h\|_{0,\Omega}.
\end{equation}
\end{Lemma}

结合 Lemma~\ref{lem:incompressible-superconvergence-mapped-gradient-identity}
至 Lemma~\ref{lem:incompressible-superconvergence-cross-terms},
得到下列 Theorem.

\begin{Theorem}\label{thm:incompressible-quadrilateral-superconvergence}
令 $\Omega$ 和 $\mathcal T_h$ 如
Lemma~\ref{lem:incompressible-superconvergence-cross-terms} 中所述,
$I_\kappa$ 是 $\kappa$ 上与 $I$ 对应的插值算子:
\[
(I_\kappa\phi)\circ F_\kappa
=I(\phi\circ F_\kappa)
\qquad\text{on each }\kappa,
\]
其中 $I\in\mathcal L(C^0(\widehat\kappa);Q_l)$
是在 $(l+1)^2$ 个 Gauss--Lobatto 求积节点上的标准插值算子.
则存在与 $h$ 无关的常数 $C>0$, 使对所有
$\phi\in H^{l+2}(\Omega)$,
\begin{equation}\label{eq:incompressible-quadrilateral-superconvergence}
\sup_{\mu_h\in\Theta_h}
\frac{|(\operatorname{curl}(\phi-I_h\phi),
\operatorname{curl}\mu_h)|}{\|\mu_h\|_{0,\Omega}}
\leq Ch^l\|\phi\|_{l+2,\Omega}.
\end{equation}
\end{Theorem}

利用上一小节的材料, 立即得到下列推论.

\begin{Corollary}\label{cor:incompressible-quadrilateral-kernel-superconvergence}
在 Theorem~\ref{thm:incompressible-quadrilateral-superconvergence}
的假设下, 存在与 $h$ 无关的常数 $C>0$, 使
\begin{equation}\label{eq:incompressible-quadrilateral-kernel-superconvergence}
\inf_{\mathbf v_h\in\widetilde V_h}|\mathbf v-\mathbf v_h|
\leq Ch^l\|\phi\|_{l+2,\Omega}
\end{equation}
对所有
$\mathbf v=(\operatorname{curl}\phi,\theta)\in\widetilde V$
且 $\phi\in H^{l+2}(\Omega)$ 成立.
\end{Corollary}
